\documentclass[11pt]{article}
\usepackage[letterpaper,margin=1in]{geometry}
\usepackage{amsmath,amssymb,lmodern,microtype}
\usepackage[T1]{fontenc}
\usepackage[hidelinks]{hyperref}
\begin{document}
\begin{center}
{\LARGE How large must a triangle be to survive?}\\[6pt]
{\large Minimum-distance exclusion under two sampling laws}\\[6pt]
Jonathan R. Landers
\end{center}

A triangle can have a long diameter and still have two vertices too close
together to satisfy a minimum-distance requirement. Enlarging the triangle
eventually resolves the pair. How large must a randomly drawn shape become?
We compare two ways of drawing a triangle. One loses a greater fraction of
shapes at every diameter above the cutoff. Its mean required diameter is
infinite, although each individual shape has a finite threshold.

\paragraph{Shapes and scale.}
Take the three vertices as the required points and let $k>0$ be the minimum
permitted distance between distinct points. The diameter $D$ is the longest
side. If $S$ is the ratio of shortest to longest side, the closest two
vertices lie $DS$ apart. The shape first satisfies the cutoff at diameter
$M=k/S$, with equality allowed. Thus the fraction lost at a given diameter
$D$ is $\Pr(M>D)=\Pr(S<k/D)$.

Under law $R$, choose an acute angle $\theta$ uniformly in $(0,\pi/2)$ and
form a right triangle with unit hypotenuse. Under law $A$, fix one side at one
and draw the other two lengths $(u,v)$ uniformly by area from
$T=\{(u,v)\in(0,1]^2:u+v>1\}$. The unit side is longest and the inequality
ensures a nondegenerate triangle. Both laws fix the same diameter before
rescaling.

\paragraph{Theorem.}
\textit{For every $D>k$, the minimum surviving diameters obey
$\Pr(M_R>D)>\Pr(M_A>D)$. Their means satisfy
$\mathbb E[M_R]=\infty$ and $\mathbb E[M_A]=4k\log 2$.}

\paragraph{Derivation of the exclusion probabilities.}
Write $t=k/D$ and $E_i(t)=\Pr(S_i<t)$ for $i=R,A$. The legs of the right
triangle are $\sin\theta$ and $\cos\theta$. For $0<t\le1/\sqrt2$, both reach
the cutoff exactly when $\arcsin t\le\theta\le\pi/2-\arcsin t$.
Dividing the excluded angular length by $\pi/2$ gives
\[
 E_R(t)=
 \begin{cases}
 (4/\pi)\arcsin t,&0<t\le1/\sqrt2,\\
 1,&t>1/\sqrt2 .
 \end{cases}
\]
At $t=1/\sqrt2$ only the isosceles right triangle survives, a single angle of
probability zero.

For law $A$, $S_A=\min(u,v)$ and $T$ has area $1/2$. When $t\le1/2$, the
corner with $u<t$ has area $\int_0^t u\,du=t^2/2$, since $1-u<v\le1$.
The corner with $v<t$ has the same area and the two are disjoint. For
$1/2<t<1$, the surviving square $[t,1]^2$ lies entirely in $T$. Dividing
by the total area gives
\[
 E_A(t)=
 \begin{cases}
 2t^2,&0<t\le1/2,\\
 1-2(1-t)^2,&1/2<t<1,\\
 1,&t\ge1 .
 \end{cases}
\]
\newpage
\paragraph{Ordering at every diameter.}
At $D=2k$ the two curves give losses of $2/3$ and $1/2$, respectively.
To see why their order persists, bound $E_A(t)$ by a straight line through
the origin. The least possible slope is the largest value of
\[
 \frac{E_A(t)}t=
 \begin{cases}
 2t,&0<t\le1/2,\\
 4-2t-1/t,&1/2<t\le1 .
 \end{cases}
\]
The first branch increases to one. The second has derivative $-2+1/t^2$,
positive below $1/\sqrt2$ and negative above. The maximum slope is
therefore $4-2\sqrt2\simeq1.172$, below $4/\pi\simeq1.273$.
Since $\arcsin t\ge t$, for $0<t\le1/\sqrt2$ we have
\[
 E_A(t)\le(4-2\sqrt2)t < (4/\pi)t\le E_R(t).
\]
For $1/\sqrt2<t<1$, the right-triangle loss is already one, while the
side-sampled law retains a square of positive area. Together these intervals
cover every $D>k$. At $D=k$ only an equilateral triangle can pass, and its
probability under law $A$ is zero. The two loss probabilities then coincide.

\paragraph{Mean required diameter.}
The same narrow shapes that determine the large-$D$ loss also determine
whether $M$ has a finite mean. For a right triangle, symmetry exchanges the
two legs at $\theta=\pi/4$. On the lower half of the angular interval,
$S_R=\sin\theta$, so
\[
 \frac{\mathbb E[M_R]}k
 =\frac4\pi\int_0^{\pi/4}\frac{d\theta}{\sin\theta}
 \ \ge\ \frac4\pi\int_0^{\pi/4}\frac{d\theta}{\theta}
 =\infty .
\]
Under law $A$, differentiating the continuous, piecewise smooth function
$E_A$ gives density $4s$ for $0<s<1/2$ and $4(1-s)$ for $1/2<s<1$.
Consequently
\begin{align*}
 \frac{\mathbb E[M_A]}k
 &=\int_0^{1/2}\frac{4s}{s}\,ds+
   \int_{1/2}^1\frac{4(1-s)}s\,ds\\
 &=2+4(\log 2-1/2)=4\log 2 .
\end{align*}
This proves both claims. \hfill$\square$

The difference comes from the room each law gives to a short side. Under
uniform angles, a leg shorter than $t$ occupies an angular interval of order
$t$ at either endpoint. Under uniform side pairs, a side of length $u<t$ also
confines its partner to an interval of width $u$. Integrating that width gives
a loss of order $t^2$. At large diameters the excluded fractions thus decay
as $D^{-1}$ and $D^{-2}$, respectively. The choice of how to sample a shape
determines whether its mean required diameter exists.

\smallskip
{\small\noindent\textit{Related work.}
The angular calculation follows J.~R. Landers,
\textit{Consequences of Minimum Distance in Mathematics and Physics},
section on surviving triangles. For side distributions obtained by sampling
directions independently, see
\href{https://www.cambridge.org/core/journals/journal-of-applied-probability/article/size-distributions-in-random-triangles/117ED3787DE9B79B1AC3862B0B6C3EAB}{\textit{Size distributions in random triangles}},
\textit{Journal of Applied Probability}.}
\end{document}
